%!TEX root = paper.tex

\section{Preliminaries}
\label{sec:prelim}

\subsubsection{Ontologies and Queries} We use standard first-order
logic  and assume familiarity with
description logics, ontology languages and theorem proving. 
A signature $\Sigma$ is a set of predicates and
 $\sig{F}$ denotes the signature of a set of formulas $F$.
It is assumed that the nullary falsehood predicate $\bot$ belongs to every $\S$. 
To capture a wide range of KR languages, we formalise
ontology axioms as  \emph{rules}: function-free sentences 
of the form
$\forall \x.[\fml(\x)
\rightarrow\exists\y.[\bigvee_{i=1}^n\fmm_i(\x,\y)]]$,
where $\fml$, $\fmm_i$ are conjunctions of distinct atoms. 
Formula $\fml$ is the rule \emph{body}  and $\exists
\y.[\bigvee_{i=1}^n\fmm_i(\x,\y)]$ is the \emph{head}. 
Universal quantification is omitted for brevity.
Rules are required to be
safe (all variables in the head occur in the body) and
we assume
w.l.o.g.\ that $\top$ (resp.\ $\bot$) does not occur in rule heads 
(resp.\ in rule bodies). A TBox
$\T$ is a finite set of rules; 
TBoxes mentioning equality $(\approx)$ are extended
with its standard axiomatisation.
A fact $\fct$ is a function-free ground atom. 
An ABox $\abox$ is a finite set of facts.
%Our view of ontologies is very general and captures
%description logics as well as most existing rule-based KR formalisms. 
A \emph{positive existential query (PEQ)} is a formula $q(\x) = \exists \y.
\fml(\x,\y)$, where $\fml$ is built from function-free atoms using only $\land$ and~$\lor$.
  

\subsubsection{Datalog}
A rule is \emph{datalog} if its head has at most one atom
and all variables are universally quantified.
A \emph{datalog program} $\prog$ is a set of datalog rules.
Given $\prog$ and an ABox $\abox$, their
\emph{materialisation} is the
set of facts entailed by $\prog \cup \abox$, which
can be computed by means of forward-chaining. A fact $\fct$ is a consequence of a datalog rule 
$r=\bigwedge_{i=1}^n\fct'_i\to\fcu$ and facts
$\fct_1,\dots,\fct_n$ if $\fct = \fcu\s$  with $\s$
a most-general unifier (MGU) of $\fct_i,\fct'_i$ for each $1 \leq i \leq n$.
A (forward-chaining) \emph{proof}  
of $\fct$ in $\prog \cup \abox$  is a pair $\rho =
(T, \lambda)$ where $T$ is a tree, $\lambda$ is a mapping from 
nodes in $T$ to facts,
and from edges in $T$ to rules in $\prog$,
such that for each node $v$
the following holds:
\begin{inparaenum}
\item $\lambda(v) = \fct$ if $v$ is the root of $T$;
\item $\lambda(v) \in
\abox$ if $v$ is a leaf;
and  
\item if $v$ has children
$w_1,\dots,w_n$ then each edge from $v$
to $w_i$ is labelled by the same rule $r$ and
$\lambda(v)$ is a consequence of $r$
and $\lambda(w_1), \dots, \lambda(w_n)$.
\end{inparaenum}
Forward-chaining is sound and complete:
a fact $\fct$ is in the materialisation 
of $\prog\cup\abox$ iff it has a proof in
$\prog \cup \abox$.
Finally, the \emph{support} of  $\fct$ is
the set of rules occurring in some proof of
$\fct$ in $\prog \cup \abox$.


\subsubsection{Inseparability Relations \& Modules}
We next recapitulate the most common 
inseparability relations studied in the literature.
We say that 
TBoxes $\T$ and 
$\T'$ are 
%
\begin{itemize}
  \item \emph{$\S$-model inseparable} ($\T\eqm_{\S}\T'$), 
if for every model $\calI$ of $\T$ (resp.\ of $\T'$) 
there exists a model $\calJ$ of $\T'$ (resp.\ of $\T$) with the same domain
s.t.\ $\A^{\calI}=\A^{\calJ}$ for each $\A\in\S$.
 \item \emph{$\S$-query inseparable} ($\T\eqq_\S\T'$) if for every
    Boolean PEQ $q$ and $\S$-ABox $\abox$ we have
    $\T\cup\abox\models q$ iff $\T'\cup\abox\models q$.
 \item \emph{$\S$-fact inseparable} ($\T\eqf_{\S}\T'$) if for every
    fact $\fct$ and ABox $\abox$ over $\S$ we have $\T\cup\abox
    \models \fct$ iff $\T'\cup \abox\models \fct$.\qedhere
 \item \emph{$\S$-implication inseparable} ($\T\eqi_{\S}\T'$) if for
 each $\varphi$ of the form $\A(\x)\rightarrow\B(\x)$ with 
 $\A,\B\in\S$,  \hbox{$\T\models
    \varphi$} iff $\T'\models
  \varphi$.
\end{itemize}

These relations are naturally ordered
from strongest to weakest:
${\eqm_{\S}}\subsetneq{\eqq_{\S}}\subsetneq{\eqf_{\S}}\subsetneq{\eqi_{\S}}$
for each
non-trivial $\Sigma$.
 

Given an inseparability relation $\equiv$ for $\S$, 
a subset $\M\subseteq \T$ is a \emph{$\equiv$-module of\/
$\T$} if $\T\equiv \M$. Furthermore, $\M$ is \emph{minimal}
if no $\M' \subsetneq \M$ is a $\equiv$-module of\/
$\T$.
